\head{第二十二版整合核查：第14节的实际范围}
\audit{V22-INT01}{实际接续点与逐段范围}{从第二十一版源码接续，完整回校已有第14节}
本轮基线为公开提交\texttt{79c1100b}的\path{current/EGA_v21}。其上一轮回校确实止于第一篇第13节；第14节已有中文译文，但尚未接受本轮法英对读。本版将该节按原稿次序重校为三个连续模块，从节标题及附记至第15节之前的最后一个交换图和说明，不把已有译文再次计为从未翻译过的内容。

\begin{longtable}{@{}p{.17\linewidth}p{.23\linewidth}p{.48\linewidth}@{}}
\toprule 模块 & 实际底本位置 & 连续内容边界\\\midrule\endhead
锥投影的几何 & 法83--84及85页上部；英53末标题、54--55 & 商约定、中心、关联空间、两个交换图、一般中心及层记号，止于说明两处$F$不同。\\
纤维与平坦性 & 法85末--88首段；英55末--57 & 从“粗略地说”开始，正则中心、爆破、14.1--14.3及模版本，止于CM条件的必要性。\\
微分性质 & 法88第二段--90的第15节之前；英57--59 & 微分性质导言、14.4--14.6、非浸入及非分歧限制、一般基底图式和函子记号。\\
\bottomrule
\end{longtable}

页码均按各完整PDF的物理页从1起数。英文第53页末只有第14节标题；法稿正文从物理83页开始。法90末、英59出现的第15节标题是下一回校点，不属于本轮新审内容。另读的相邻页及法21页只服务接缝或引用核对，不扩充本轮完成范围。逐段映射、实际源页图和文件散列保存在各模块证据目录中。

\audit{V22-INT02}{原错、转写错与补证的分界}{正文采用可证明的版本，脚注保存原断言及订正依据}
\lead{可直接据法稿恢复的读法。}
包括\emph{monomorphisme}、闭子预概形、13.4(c)、一般截面下标、14.1(b)的像维数与“非空”纤维、被英文漏掉的反证条件、正确的纤维积、模版本的写作指示，以及第14节末图式中的$S$和波浪号。英文编者或译者的插注仍标明身份；已可辨的法文不再沿用英文的“不可辨”标记。

\lead{14.1与14.3的假设确实不足。}
爆破$\operatorname{Bl}_p(\mathbb P^3)\to\mathbb P^3$具有零维一般纤维，但投影中仍可出现包含例外$\mathbb P^2$的纤维。因此14.1(a)不能只控制一般纤维；本版明列各不可约分支的全部非空几何纤维纯$d$维这一充分条件。可约CM源还可能有低维像分支：一个平面与其外孤立点的并给出14.3不平坦的反例。正文另加各分支像维数至少$m$，报告分别证明纯维数、总空间CM、平坦性及CM纤维，不能仅凭最大纤维维数相等推出平坦。

\lead{14.4的边界及手添强断言。}
14.4(a)关于有限点集的根态射结论需有$m\geq1$；两个点投向$\mathbb P^0$已经构成反例。另一断言必须要求$Y$与指定点集$T$不交，否则$Y=T$即失败。这些限制没有误加到正确的14.4(b)上。法稿把非分歧性加强到$p^{-1}p(T)$的手添句可以辨读，但异面直线及另一条直线的平方理想加厚构成反例。因此正文只保留在$T$上非分歧的已证范围，原加强句及反例移入明确标识的脚注和报告。

\lead{14.5不能在任意闭点处平展。}
特征零的整曲线$\operatorname{Spec}k[t^3,t^4,t^5]$在原点的嵌维为3，而其平面投影像的切空间维数至多为2；该点处不可能非分歧，更不可能平展。本版在有限型、几何约化范围内，对原光滑闭点给出正确充分版本。补证使用整个射影闭包，排除边界点落入同一纤维；随后分别核验唯一几何纤维、剩余域相等、非分歧性及Nakayama引理，才得到局部同构和双有理性。对每个原闭点分别选非空好开并取泛点，不把无限多个开集的交冒称为含有有理点的非空开。

\lead{补证不替换原稿的写作身份。}
任意拟凝聚模只在已核验的Koszul意义下使用一般中心的正则性；普通正则列及CM模结论明列有限性。对正则中心的爆破图表、一般纤维维数和CM平坦性的独立证明，与原作者简略提示分别排印。第14节导言关于普通奇点、唯一性及有理性的研究设想仍按草稿措辞保留，未被写成已经完成的一般理论。

\lead{仍未解决的内容。}
层记号括号内的一处手写字母不能可靠判读；正文采用依赖中心的记号并说明编辑选择。原“第21页上方、第5节之前”引用尚未定位。原“可分”术语在全套手稿中的全部使用范围也未在本轮作统一考证；正面补证明确自己的几何约化条件。原8.7(c)、8.8、9.5的缺文及13.4(d)任意基底推广等旧限制继续保留，不因本轮更新而消失。

\audit{V22-INT03}{源码保留、构建与交接}{文件一致性与数学审定各有边界}
第二十一版的194个正文、续篇、书末及报告TeX文件作为继承对象原字节保留。原\path{continuations/v_dot_late.tex}共49,224字节，其中第14节前缀34,610字节、第15节后缀14,614字节；二者拼接还原原文件。三份新模块在活动输入链中替换原第14节，旧全文及47条原译注保存在历史证据中，第15节精确后缀只作为继承译文接入。Dot69--108的40个活动模块仍各输入一次；旧工程、后续Dot候选稿和累计核查报告继续保存。

本版新增41条第14节译注，三份专题报告共19项核查，另有本节三项整合记录。这些计数不表示原稿新增了相同数量的定理或错误。新版报告把历版意见完整放在后面，历史“本轮”只指各自版次。

使用真正的\texttt{ctexbook}默认字体、字号、版心和分页，沿用原导言文件；公式和交换图由TeX排印。工程对当前模块与独立单文件分别编译，并记录最终日志、全页文字与尺寸及72dpi像素一致性、明确列出的成品页面目视检查。实际页数和看图范围以随工程保存的检查记录为准，不能把抽样页面的目视通过写成全书每页人工数学审核。

下一回校点为第一篇第15节，法文物理90页末、英文59页，现有译文在\path{continuations/v22_section15_inherited.tex}。正式第IV章主体仍从第四卷第一分册物理221页提要及后续章首接续，属于另一个未完成分支。全套书籍尚未完结；历史中7个早期原ZIP容器不可恢复的事实仍保留，普通成员文件的留存不等于原容器已经找回。
